\section{Introduction}


In recent years, training locomotion policies in simulation using reinforcement learning (RL) then transferring them to real robots has proven highly effective in mastering agile skills~\citep{lee2020learning, bellegarda2020robust, fu2021minimizing,margolis2021learning, rudin2022learning, bellegarda2022robust,li2023robust, he2024agile,inoue2024body, ma2024dreureka}.
Notably,~\citep{caluwaerts2023barkour, zhuang2023robot, cheng2023parkour, hoeller2024anymal} have successfully demonstrated a range of dynamic skills akin to the athletic maneuvers seen in parkour, including walking, climbing, jumping, and crawling, which necessitates tight coordination between vision and control.
In this work, we are interested in developing parkour skills from onboard depth inputs for Solo-12, a lightweight, open-hardware quadruped robot \cite{grimminger2020open}, using this sim-to-real approach. 
Solo is a lighter prototype than other quadruped robots experimented in parkour projects (e.g. Atlas, Go-1, or Anymal). Safe visual locomotion with Solo is then at its highest stake, to prevent exceeding the limited torque range of its motors and breaking the 3D-printed plastic shells and exposed electronic components in unexpected impacts or falls. 

Yet, training visual locomotion policies presents significant challenges due to the limited field of view and dynamic, lag-prone visual perception during dynamic movements.
Previous methods for end-to-end control from pixels~\citep{agarwal2023legged, zhuang2023robot, cheng2023parkour} typically involve a two-stage process,
where a policy is first trained with privileged information about the robot's surroundings, then the learned behaviors are distilled into a visual policy using imitation learning~\citep{ross2011reduction}.
However, such an approach relies on the unrealistic assumption that privileged information can be fully reconstructed from a history of depth images.
Indeed, such privileges may reveal information obscured behind obstacles, outside the field of view of the egocentric camera, or simply missing from the depth image stream due to lags.
This gap between the privileged information and the actual visual inputs may result in behaviors that a vision-based policy cannot replicate accurately.
In that case, an optimal vision policy would differ from the one using privileged information.
It would, for example, try to gather more information before crossing an obstacle to handle the occlusion.
Ideally, training RL directly from pixels would ensure the policy learns behaviors effectively based on its actual visual sensors, but this method is impractical due to the high computational costs associated with generating depth images in simulations~\citep{agarwal2023legged}.

To address these challenges, we propose \textit{SoloParkour}, a novel approach for training parkour locomotion policies using depth inputs.
Our method frames parkour as a constrained reinforcement learning problem~\citep{lee2023evaluation, kim2024not, chane2024cat}.
This framework allows the RL agent to explore the full capabilities of the robot while explicitly preventing unsafe actions, ensuring that the agent can aggressively optimize performance without compromising the safety of the robot.
We propose a novel method to train the visual locomotion policy end-to-end from pixels using a sample-efficient off-policy RL algorithm.
Building upon recent advancements in accelerating RL using demonstrations from previous controllers~\citep{ball2023efficient, smith2023learning}, our algorithm not only learns from its own experiences but also leverages experience generated by the privileged policy.
This approach enables the RL agent to rapidly acquire agile skills and adapt behaviors from the privileged experience to suit the specific limitations of its visual sensors, bypassing the computational burden associated with RL from pixels while still ensuring the development of agile and safe legged locomotion skills.

We demonstrate the effectiveness of our approach in simulation.
We then directly deploy the learned policy on a real Solo-12 robot and demonstrate the effective acquisition of parkour skills such as crawling, leaping, and climbing obstacles from depth image (see Figure~\ref{fig:teaser}).
Our approach pushes the robot to its limits: it manages to clear obstacles 1.5 times higher than its height despite the robot being significantly less powerful than the ones typically used in parkour experiments.

In summary, our contributions are as follows:
\begin{itemize}
\item we cast parkour learning from depth images as a constrained RL problem,
\item we introduce a computationally efficient RL algorithm to train end-to-end visual locomotion policies with significant improvement over methods based on distillation
\item and we validate the effectiveness of our approach in simulation and on a real Solo-12 robot to perform parkour skills, outperforming the best movements ever generated with this robot and reaching performances comparable to recent parkour achievements despite a more limited actuation range.
\end{itemize}
